\section{Environment Isolation and Reproducibility}
\label{sec:reproducibility}

A defining characteristic of scientific software is reproducibility.
Deep learning experiments are notoriously sensitive to variations in the software stack, where differences in library versions or CUDA drivers can lead to non-deterministic numerical divergences.
To mitigate this, the entire training pipeline is containerized using \textbf{Docker}.

\subsection{Containerization Strategy}
The experimental environment is strictly defined via a \texttt{Dockerfile} that builds upon the official NVIDIA PyTorch image (\texttt{nvcr.io/nvidia/pytorch:24.06-py3})~\cite{nvidia_pytorch_docker}.
This base image provides a validated stack of low-level dependencies, including the CUDA Toolkit, cuDNN primitives, and NCCL communication libraries, optimized for the underlying GPU hardware.

By encapsulating the project dependencies (including specific versions of \texttt{PyTorch 2.4.0}, \texttt{Optuna}, and scientific libraries like \texttt{scikit-image}) within an immutable image, we guarantee that the code executes in an identical software context whether it is running on a local Windows development workstation or on a distributed Linux HPC cluster.
This approach elevates the project from a set of experimental scripts to deployable scientific software, ensuring that any researcher can replicate the exact experimental conditions by simply building the container.

